\documentclass{ieeeaccess}

\usepackage{cite}
\usepackage{amsmath,amssymb,amsfonts}
\usepackage{algorithmic}
\usepackage{graphicx}
\usepackage{textcomp}
\usepackage{xcolor}
\usepackage{booktabs}
\usepackage{multirow}

\def\BibTeX{{\rm B\kern-.05em{\sc i\kern-.025em b}\kern-.08em
    T\kern-.1667em\lower.7ex\hbox{E}\kern-.125emX}}

\begin{document}

\history{Date of publication xxxx 00, 0000, date of current version xxxx 00, 0000.}
\doi{10.1109/ACCESS.2026.DOI}

\title{Risk-Aware Soft Actor-Critic for Dynamic Obstacle Avoidance of a UR5 Manipulator}

\author{\uppercase{First Author}\authorrefmark{1},
\uppercase{Second Author}\authorrefmark{1}, and
\uppercase{Corresponding Author}\authorrefmark{1}}

\address[1]{School/College, Shanghai University, Shanghai, China}

\corresp{Corresponding author: Corresponding Author (e-mail: author@shu.edu.cn).}

\begin{abstract}
Dynamic obstacle avoidance is a critical capability for robotic manipulators operating in unstructured and human-shared environments.
This paper proposes a risk-aware soft actor-critic framework for dynamic obstacle avoidance of a UR5 manipulator.
The method represents the manipulator links with capsule-based geometric primitives and incorporates link-level obstacle proximity into the reinforcement learning reward.
By penalizing unsafe configurations and collisions during training, the learned policy is encouraged to reach target positions while maintaining safe distances from moving obstacles.
Simulation experiments under randomized targets and dynamic obstacles are conducted to evaluate success rate, collision rate, minimum link-obstacle distance, trajectory efficiency, and policy robustness.
The proposed method is compared with standard SAC and ablated variants to verify the contribution of link-level risk modeling.
Experimental results show that the risk-aware formulation improves obstacle avoidance safety while preserving task performance.
\end{abstract}

\begin{keywords}
Robot motion planning, reinforcement learning, soft actor-critic, robotic manipulators, obstacle avoidance, risk-aware control, collision avoidance, dynamic environments, UR5 manipulator.
\end{keywords}

\titlepgskip=-15pt
\maketitle

\section{Introduction}
\label{sec:introduction}

Robotic manipulators are increasingly expected to operate in dynamic and partially structured environments.
In such scenarios, obstacle avoidance should consider not only the end-effector but also the whole body of the robot, since collisions may occur along any link of the manipulator.

Write this section around four questions:
\begin{itemize}
    \item Why is dynamic obstacle avoidance important for robotic manipulators?
    \item Why are traditional planning or control methods insufficient in the target setting?
    \item Why can reinforcement learning help, and why does ordinary reinforcement learning still have safety limitations?
    \item What are the main contributions of the proposed risk-aware SAC method?
\end{itemize}

The main contributions of this paper are summarized as follows:
\begin{itemize}
    \item A risk-aware reinforcement learning framework is proposed for dynamic obstacle avoidance of a UR5 manipulator.
    \item A link-level capsule representation is introduced to evaluate proximity risk between the manipulator body and dynamic obstacles.
    \item The risk information is integrated into the SAC training objective through a safety-aware reward design.
    \item Comparative and ablation experiments are conducted to evaluate task success, collision avoidance, and robustness.
\end{itemize}

\section{Related Work}
\label{sec:related_work}

\subsection{Obstacle Avoidance for Robotic Manipulators}
Discuss sampling-based planning, optimization-based trajectory generation, artificial potential fields, model predictive control, and their limitations in dynamic environments.

\subsection{Reinforcement Learning for Manipulator Control}
Discuss DDPG, TD3, PPO, SAC, and continuous-control robotic manipulation.

\subsection{Safe and Risk-Aware Reinforcement Learning}
Discuss constrained RL, risk-sensitive RL, collision penalties, control barrier functions, and safety-aware reward shaping.

\section{Problem Formulation}
\label{sec:problem}

Define the dynamic obstacle avoidance task as a Markov decision process.
Let $s_t$ denote the state, $a_t$ denote the action, and $r_t$ denote the reward at time step $t$.
The objective is to learn a policy $\pi(a_t|s_t)$ that drives the UR5 end-effector toward a target while avoiding collisions between any robot link and moving obstacles.

\subsection{State and Action Spaces}
Describe joint positions, joint velocities, end-effector pose, target position, obstacle position, obstacle velocity, and optional risk features.

\subsection{Collision and Safety Definitions}
Define the minimum distance between link capsules and obstacles:
\begin{equation}
    d_{\min}(s_t) = \min_{i,j} d(C_i(q_t), O_j(t)),
\end{equation}
where $C_i(q_t)$ is the capsule of the $i$-th link under joint configuration $q_t$, and $O_j(t)$ is the $j$-th obstacle.

\section{Method}
\label{sec:method}

\subsection{Capsule-Based Link-Level Risk Modeling}
Describe how each UR5 link is approximated as a capsule and how the link-obstacle distance is computed.

\subsection{Risk-Aware Reward Design}
An example reward decomposition is:
\begin{equation}
    r_t = r_{\mathrm{goal}} + r_{\mathrm{progress}} + r_{\mathrm{safety}} + r_{\mathrm{collision}} + r_{\mathrm{action}} + r_{\mathrm{time}}.
\end{equation}

Explain every term and connect it to the implementation.

\subsection{Soft Actor-Critic Training}
Describe actor, critic, entropy regularization, replay buffer, target network update, and the training loop.

\subsection{Predictive Risk Extension}
If the final experiment uses predictive risk, describe short-horizon obstacle prediction and the corresponding future-distance penalty here.
If not used, remove this subsection.

\section{Experiments}
\label{sec:experiments}

\subsection{Experimental Setup}
Describe simulator, UR5 model, random target generation, dynamic obstacle generation, training steps, evaluation episodes, random seeds, and hardware.

\subsection{Baselines}
Suggested baselines:
\begin{itemize}
    \item Standard SAC without explicit risk penalty.
    \item SAC with end-effector-only obstacle distance penalty.
    \item The proposed link-level risk-aware SAC.
\end{itemize}

\subsection{Evaluation Metrics}
Use metrics such as success rate, collision rate, average return, final distance to target, minimum link-obstacle distance, episode length, path length, and action smoothness.

\subsection{Main Results}
Insert the main comparison table and discuss safety-performance tradeoffs.

\subsection{Ablation Study}
Discuss the effects of risk weight, safety threshold, link-level distance, and predictive risk.

\subsection{Robustness Analysis}
Evaluate the method under different obstacle speeds, obstacle counts, and target distributions.

\section{Discussion}
\label{sec:discussion}

Discuss practical implications, limitations, and sim-to-real considerations.
Be explicit if the current work is simulation-only.

\section{Conclusion}
\label{sec:conclusion}

This paper presented a risk-aware SAC framework for dynamic obstacle avoidance of a UR5 manipulator.
By incorporating link-level geometric risk into reinforcement learning, the proposed method improves collision avoidance behavior in dynamic environments.
Future work will investigate real-robot validation, richer perception inputs, and formal safety constraints.

\section*{Acknowledgment}

Add funding information and any required AI-use disclosure here.

\begin{thebibliography}{00}

\bibitem{haarnoja2018sac}
T. Haarnoja, A. Zhou, P. Abbeel, and S. Levine, ``Soft actor-critic: Off-policy maximum entropy deep reinforcement learning with a stochastic actor,'' in \emph{Proc. Int. Conf. Mach. Learn.}, 2018, pp. 1861--1870.

\bibitem{khatib1986real}
O. Khatib, ``Real-time obstacle avoidance for manipulators and mobile robots,'' \emph{Int. J. Robot. Res.}, vol. 5, no. 1, pp. 90--98, 1986.

\end{thebibliography}

\begin{IEEEbiography}{First Author}
Biography text here. Include education, current affiliation, research interests, and major achievements.
\end{IEEEbiography}

\begin{IEEEbiography}{Second Author}
Biography text here.
\end{IEEEbiography}

\EOD

\end{document}
